% §\,3. Функция
% Зорич, Математический анализ, Том I
% Книжные страницы 10--23 (PDF-страницы 22--35)

\section*{\S\,3. Функция}

\subsection*{1. Понятие функции (отображения)}

Перейдем теперь к описанию фундаментального не только для математики понятия функциональной зависимости.

Пусть $X$ и $Y$ --- какие-то множества.

Говорят, что имеется \textit{функция}, определенная на $X$ со значениями в $Y$,
если в силу некоторого закона $f$ каждому элементу $x \in X$ соответствует элемент
$y \in Y$.

В этом случае множество $X$ называется \textit{областью определения} функции;
символ $x$ его общего элемента --- \textit{аргументом функции} или \textit{независимой переменной};
соответствующий конкретному значению $x_0 \in X$ аргумента $x$ элемент
$y_0 \in Y$ называют \textit{значением функции на элементе} $x_0$ или значением
функции при значении аргумента $x = x_0$ и обозначают через $f(x_0)$. При изменении
аргумента $x \in X$ значения $y = f(x) \in Y$, вообще говоря, меняются в
зависимости от значений $x$. По этой причине величину $y = f(x)$ часто называют
\textit{зависимой переменной}.

Множество
\[
  f(X) := \{ y \in Y \mid \exists x\; ((x \in X) \land (y = f(x))) \}
\]
всех значений функции, которые она принимает на элементах множества $X$,
будем называть \textit{множеством значений} или \textit{областью значений функции}.

В зависимости от природы множеств $X$, $Y$ термин <<функция>> в различных
отделах математики имеет ряд полезных синонимов: \textit{отображение}, \textit{преобразование}, \textit{морфизм}, \textit{оператор}, \textit{функционал}. Отображение --- наиболее распространенный из них, и мы его тоже часто будем употреблять.

Для функции (отображения) приняты следующие обозначения:
\[
  f \colon X \to Y, \quad X \xrightarrow{f} Y.
\]

Когда из контекста ясно, каковы область определения и область значений
функции, используют также обозначения $x \mapsto f(x)$ или $y = f(x)$, а чаще обозначают функцию вообще одним лишь символом $f$.

Две функции $f_1$, $f_2$ считаются \textit{совпадающими} или \textit{равными}, если они имеют одну и ту же область определения $X$ и на любом элементе $x \in X$ значения
$f_1(x)$, $f_2(x)$ этих функций совпадают. В этом случае пишут $f_1 = f_2$.

Если $A \subset X$, а $f \colon X \to Y$ --- некоторая функция, то через $f|A$ или $f|_A$ обозначают функцию $\varphi \colon A \to Y$, совпадающую с $f$ на множестве $A$, т.\,е. $f|_A(x) =
{} = \varphi(x) = f(x)$, если $x \in A$. Функция $f|_A$ называется \textit{сужением} или \textit{ограничением} функции $f$ на множество $A$, а функция $f \colon X \to Y$ по отношению к функции
$\varphi = f|_A \colon A \to Y$ называется \textit{распространением} или \textit{продолжением} функции $\varphi$
на множество $X$.

Мы видим, что иногда приходится рассматривать функцию $\varphi \colon A \to Y$,
определенную на подмножестве $A$ некоторого множества $X$, причем область
значений $\varphi(A)$ функции $\varphi$ тоже может оказаться не совпадающим с $Y$
подмножеством множества $Y$. В связи с этим для обозначения любого множества $X$, содержащего область определения функции, иногда используется
термин \textit{область отправления функции}, а любое множество $Y$, содержащее
область значений функции, называют тогда \textit{областью её прибытия}.

Итак, задание функции (отображения) предполагает указание тройки
$(X, f, Y)$, где

$X$ --- отображаемое множество, или область определения функции;

$Y$ --- множество, в которое идет отображение, или область прибытия функции;

$f$ --- закон, по которому каждому элементу $x \in X$ сопоставляется определенный элемент $y \in Y$.

Наблюдаемая здесь несимметричность между $X$ и $Y$ отражает то, что
отображение идет именно из $X$ в $Y$.

Рассмотрим некоторые примеры функций.

\medskip

\textbf{Пример~1.} Формулы $l = 2\pi r$ и $V = \dfrac{4}{3}\,\pi r^3$ устанавливают функциональную
зависимость длины окружности $l$ и объема шара $V$ от радиуса $r$. По смыслу
каждая из этих формул задает свою функцию $f \colon \mathbb{R}_+ \to \mathbb{R}_+$, определенную на
множестве $\mathbb{R}_+$ положительных действительных чисел со значениями в том
же множестве $\mathbb{R}_+$.

\textbf{Пример~2.} Пусть $X$ --- множество инерциальных систем координат, а
$c \colon X \to \mathbb{R}$ --- функция, состоящая в том, что каждой инерциальной системе
координат $x \in X$ сопоставляется измеренное относительно нее значение $c(x)$
скорости света в вакууме. Функция $c \colon X \to \mathbb{R}$ постоянна, т.\,е. при любом $x \in X$
она имеет одно и то же значение $c$ (это фундаментальный экспериментальный факт).

\textbf{Пример~3.} Отображение $G \colon \mathbb{R}^2 \to \mathbb{R}^2$ (прямого произведения $\mathbb{R}^2 = \mathbb{R} \times
{} \mathbb{R} = \mathbb{R}_t \times \mathbb{R}_x$ оси времени $\mathbb{R}_t$ и пространственной оси $\mathbb{R}_x$) на себя же, задаваемое формулами
\begin{align*}
  x' &= x - vt, \\
  t' &= t,
\end{align*}
есть классическое преобразование Галилея для перехода от одной инерциальной системы координат $(x, t)$ к другой --- $(x', t')$, движущейся относительно первой со скоростью $v$.

Той же цели служит отображение $L \colon \mathbb{R}^2 \to \mathbb{R}^2$, задаваемое соотношениями
\begin{align*}
  x' &= \frac{x - vt}{\sqrt{1 - \left(\dfrac{v}{c}\right)^2}}, \\[6pt]
  t' &= \frac{t - \dfrac{v}{c^2}\,x}{\sqrt{1 - \left(\dfrac{v}{c}\right)^2}}.
\end{align*}

Это --- известное (одномерное) \textit{преобразование Лоренца}\footnote{Г.\,А.\,Лоренц (1853--1928) --- выдающийся голландский физик-теоретик. Указанные
преобразования Пуанкаре назвал в честь Лоренца, стимулировавшего исследование симметрий уравнений Максвелла. Они существенно использованы Эйнштейном в сформулированной им в 1905 году специальной теории относительности.}, играющее фундаментальную роль в специальной теории относительности; $c$ --- скорость
света.

\textbf{Пример~4.} \textit{Проектирование} $\operatorname{pr}_1 \colon X_1 \times X_2 \to X_1$, задаваемое соответствием $X_1 \times X_2 \ni (x_1, x_2) \xmapsto{\operatorname{pr}_1} x_1 \in X_1$, очевидно, является функцией. Аналогичным
образом определяется вторая проекция $\operatorname{pr}_2 \colon X_1 \times X_2 \to X_2$.

\textbf{Пример~5.} Пусть $\mathscr{P}(M)$ --- множество всех подмножеств множества $M$.
Каждому множеству $A \in \mathscr{P}(M)$ поставим в соответствие множество $C_M A \in
{} \in \mathscr{P}(M)$, т.\,е. дополнение к $A$ в $M$. Тогда получим отображение $C_M \colon \mathscr{P}(M) \to
{} \to \mathscr{P}(M)$ множества $\mathscr{P}(M)$ в себя.

\textbf{Пример~6.} Пусть $E \subset M$. Вещественнозначную функцию $\chi_E \colon M \to \mathbb{R}$,
определенную на множестве $M$ условиями $(\chi_E(x) = 1$, если $x \in E) \land (\chi_E(x) =
{} = 0$, если $x \in C_M E)$, называют \textit{характеристической функцией множества}~$E$.

\textbf{Пример~7.} Пусть $M(X; Y)$ --- множество отображений множества $X$ в
множество $Y$, а $x_0$ --- фиксированный элемент из $X$. Любой функции $f \in
{} \in M(X; Y)$ поставим в соответствие ее значение $f(x_0) \in Y$ на элементе $x_0$.
Этим определяется функция $F \colon M(X; Y) \to Y$. В частности, если $Y = \mathbb{R}$, т.\,е.
если $Y$ есть множество действительных чисел, то каждой функции $f \colon X \to \mathbb{R}$
функция $F \colon M(X; \mathbb{R}) \to \mathbb{R}$ ставит в соответствие число $F(f) = f(x_0)$. Таким
образом, $F$ есть функция, определенная на функциях. Для удобства такие
функции называют \textit{функционалами}.

\textbf{Пример~8.} Пусть $\Gamma$ --- множество кривых, лежащих на поверхности (например, земной) и соединяющих две ее фиксированные точки. Каждой кривой $\gamma \in \Gamma$ можно сопоставить ее длину. Тогда мы получим функцию $F \colon \Gamma \to \mathbb{R}$,
которую часто приходится рассматривать с целью отыскания кратчайшей
линии или, как говорят, \textit{геодезической} линии между данными точками на
поверхности.

\textbf{Пример~9.} Рассмотрим множество $M(\mathbb{R}; \mathbb{R})$ всех вещественнозначных
функций, определенных на всей числовой оси $\mathbb{R}$. Фиксировав число $a \in \mathbb{R}$,
каждой функции $f \in M(\mathbb{R}; \mathbb{R})$ поставим в соответствие функцию $f_a \in M(\mathbb{R}; \mathbb{R})$,
связанную с ней соотношением $f_a(x) = f(x + a)$. Функцию $f_a(x)$ обычно
называют \textit{сдвигом на $a$} функции $f(x)$. Возникающее при этом отображение
$A \colon M(\mathbb{R}; \mathbb{R}) \to M(\mathbb{R}; \mathbb{R})$ называется \textit{оператором сдвига}. Итак, оператор $A$
определен на функциях и значениями его также являются функции: $f_a = A(f)$.

Рассмотренный пример мог бы показаться искусственным, если бы мы
на каждом шагу не видели реальные операторы. Так, любой радиоприемник есть оператор $f \xmapsto{F} \hat{f}$, преобразующий электромагнитные сигналы $f$ в
звуковые $\hat{f}$; любой из наших органов чувств является оператором (преобразователем) со своими областью определения и областью значений.

\textbf{Пример~10.} Положение частицы в пространстве определяется упорядоченной тройкой чисел $(x, y, z)$, называемой ее координатами в пространстве. Множество всех таких упорядоченных троек можно себе мыслить как
прямое произведение $\mathbb{R} \times \mathbb{R} \times \mathbb{R} = \mathbb{R}^3$ трех числовых осей $\mathbb{R}$.

При движении в каждый момент времени $t$ частица находится в некоторой точке пространства $\mathbb{R}^3$ с координатами $(x(t), y(t), z(t))$. Таким образом,
движение частицы можно интерпретировать как отображение $\gamma \colon \mathbb{R} \to \mathbb{R}^3$, где
$\mathbb{R}$ --- ось времени, а $\mathbb{R}^3$ --- трехмерное пространство.

Если система состоит из $n$ частиц, то ее конфигурация задается положением каждой из $n$ частиц, т.\,е. упорядоченным набором $(x_1, y_1, z_1; x_2, y_2, z_2; \ldots;$
$x_n, y_n, z_n)$ из $3n$ чисел. Множество всех таких наборов называется \textit{конфигурационным пространством системы $n$ частиц}. Следовательно, конфигурационное пространство системы $n$ частиц можно интерпретировать как прямое
произведение $\mathbb{R}^3 \times \mathbb{R}^3 \times \ldots \times \mathbb{R}^3 = \mathbb{R}^{3n}$ $n$ экземпляров пространства $\mathbb{R}^3$.

Движению системы из $n$ частиц отвечает отображение $\gamma \colon \mathbb{R} \to \mathbb{R}^{3n}$ оси
времени в конфигурационное пространство системы.

\textbf{Пример~11.} Потенциальная энергия $U$ механической системы связана с
взаимным расположением частиц системы, т.\,е. определяется конфигурацией, которую имеет система. Пусть $Q$ --- множество реально возможных конфигураций системы. Это некоторое подмножество конфигурационного пространства системы. Каждому положению $q \in Q$ отвечает некоторое значение
$U(q)$ потенциальной энергии системы. Таким образом, потенциальная энергия есть функция $U \colon Q \to \mathbb{R}$, определенная на подмножестве $Q$ конфигурационного пространства со значениями в области $\mathbb{R}$ действительных чисел.

\textbf{Пример~12.} Кинетическая энергия $K$ системы $n$ материальных частиц
зависит от их скоростей. Полная механическая энергия системы $E = K + U$,
т.\,е. сумма кинетической и потенциальной энергий, зависит, таким образом,
как от конфигурации $q$ системы, так и от набора $v$ скоростей ее частиц. Как
и конфигурация $q$ частиц в пространстве, набор $v$, состоящий из $n$ трехмерных векторов, может быть задан упорядоченным набором из $3n$ чисел.
Упорядоченные пары $(q, v)$, отвечающие состояниям нашей системы, образуют подмножество $\Phi$ в прямом произведении $\mathbb{R}^{3n} \times \mathbb{R}^{3n} = \mathbb{R}^{6n}$, называемом
\textit{фазовым пространством системы $n$ частиц} (в отличие от конфигурационного пространства $\mathbb{R}^{3n}$).

Полная энергия системы является, таким образом, функцией $E \colon \Phi \to \mathbb{R}$,
определенной на подмножестве $\Phi$ фазового пространства $\mathbb{R}^{6n}$ и принимающей значения в области $\mathbb{R}$ действительных чисел.

В частности, если система изолирована, т.\,е. на нее не действуют внешние силы, то по закону сохранения энергии в любой точке множества $\Phi$
состояний системы функция $E$ будет иметь одно и то же значение $E_0 \in \mathbb{R}$.

\subsection*{2. Простейшая классификация отображений}

Когда функцию $f \colon X \to Y$ называют отображением, значение $f(x) \in Y$, которое она принимает на
элементе $x \in X$, обычно называют \textit{образом} элемента $x$.

\textit{Образом множества} $A \subset X$ при отображении $f \colon X \to Y$ называют множество
\[
  f(A) := \{ y \in Y \mid \exists x\; ((x \in A) \land (y = f(x))) \}
\]
тех элементов $Y$, которые являются образами элементов множества $A$.

Множество
\[
  f^{-1}(B) := \{ x \in X \mid f(x) \in B \}
\]
тех элементов $X$, образы которых содержатся в $B$, называют \textit{прообразом} (или
\textit{полным прообразом}) множества $B \subset Y$ (рис.\,6).

% Рис. 6 — диаграмма с X, Y, A, f(A), f^{-1}(B), B

Про отображение $f \colon X \to Y$ говорят, что оно

\textit{сюръективно} (или есть отображение $X$ \textit{на} $Y$), если $f(X) = Y$;

\textit{инъективно} (или есть \textit{вложение}, \textit{инъекция}), если для любых элементов $x_1$,
$x_2$ множества $X$
\[
  (f(x_1) = f(x_2)) \Rightarrow (x_1 = x_2),
\]
т.\,е. различные элементы имеют различные образы;

\textit{биективно} (или \textit{взаимно однозначно}), если оно сюръективно и инъективно одновременно.

Если отображение $f \colon X \to Y$ биективно, т.\,е. является взаимно однозначным соответствием между элементами множеств $X$ и $Y$, то естественно возникает отображение
\[
  f^{-1} \colon Y \to X,
\]
которое определяется следующим образом: если $f(x) = y$, то $f^{-1}(y) = x$, т.\,е.
элементу $y \in Y$ ставится в соответствие тот элемент $x \in X$, образом которого при отображении $f$ является $y$. В силу сюръективности $f$ такой элемент
$x \in X$ найдется, а ввиду инъективности $f$ он единственный. Таким образом,
отображение $f^{-1}$ определено корректно. Это отображение называют \textit{обратным} по отношению к исходному отображению $f$.

Из построения обратного отображения видно, что $f^{-1} \colon Y \to X$ само является биективным и что обратное к нему отображение $(f^{-1})^{-1} \colon X \to Y$ совпадает с $f \colon X \to Y$.

Таким образом, свойство двух отображений быть обратными является
взаимным: если $f^{-1}$ --- обратное для $f$, то, в свою очередь, $f$ --- обратное
для $f^{-1}$.

Заметим, что символ $f^{-1}(B)$ прообраза множества $B \subset Y$ ассоциируется с
символом $f^{-1}$ обратной функции, однако следует иметь в виду, что прообраз
множества определен для любого отображения $f \colon X \to Y$, даже если оно не
является биективным и, следовательно, не имеет обратного.

\subsection*{3. Композиция функций и взаимно обратные отображения}

Богатым
источником новых функций, с одной стороны, и способом расчленения сложных функций на более простые --- с другой, является операция композиции
отображений.

Если отображения $f \colon X \to Y$ и $g \colon Y \to Z$ таковы, что одно из них (в нашем
случае $g$) определено на множестве значений другого ($f$), то можно построить новое отображение
\[
  g \circ f \colon X \to Z,
\]
значения которого на элементах множества $X$ определяются формулой
\[
  (g \circ f)(x) := g(f(x)).
\]

Построенное составное отображение $g \circ f$ называют \textit{композицией} отображения $f$ и отображения $g$ (в таком порядке!).

% Рис. 7 — диаграмма композиции f и g

С композицией отображений вы уже неоднократно встречались как в геометрии, рассматривая композицию движений плоскости или пространства,
так и в алгебре при исследовании <<сложных>> функций, полученных композицией простейших элементарных функций.

Операцию композиции иногда приходится проводить несколько раз подряд, и в этой связи полезно отметить, что она ассоциативна, т.\,е.
\[
  h \circ (g \circ f) = (h \circ g) \circ f.
\]

$\blacktriangleleft$ Действительно,
\[
  h \circ (g \circ f)(x) = h((g \circ f)(x)) = h(g(f(x))) = (h \circ g)(f(x)) = ((h \circ g) \circ f)(x). \quad \blacktriangleright
\]

Это обстоятельство, как и в случае сложения или умножения нескольких
чисел, позволяет опускать скобки, предписывающие порядок спаривания.

Если в композиции $f_n \circ \ldots \circ f_1$ все члены одинаковы и равны $f$, то ее обозначают коротко $f^n$.

Хорошо известно, например, что корень квадратный из положительного
числа $a$ можно вычислить последовательными приближениями по формуле
\[
  x_{n+1} = \frac{1}{2}\left(x_n + \frac{a}{x_n}\right),
\]
начиная с любого начального приближения $x_0 > 0$. Это не что иное, как последовательное вычисление $f^n(x_0)$, где $f(x) = \frac{1}{2}\left(x + \frac{a}{x}\right)$. Такая процедура,
когда вычисленное на предыдущем шаге значение функции на следующем
шаге становится ее аргументом, называется \textit{итерационным процессом}. Итерационные процессы широко используются в математике.

Отметим также, что даже в том случае, когда обе композиции $g \circ f$ и $f \circ g$
определены, вообще говоря,
\[
  g \circ f \neq f \circ g.
\]

Действительно, возьмем, например, двухэлементное множество $\{a, b\}$ и
отображения $f \colon \{a, b\} \to a$, $g \colon \{a, b\} \to b$. Тогда, очевидно, $g \circ f \colon \{a, b\} \to b$,
в то время как $f \circ g \colon \{a, b\} \to a$.

Отображение $f \colon X \to X$, сопоставляющее каждому элементу множества $X$
его самого, т.\,е. $x \xmapsto{f} x$, будем обозначать через $e_X$ и называть \textit{тождественным отображением множества}~$X$.

\medskip

\textbf{Лемма.}
\[
  (g \circ f = e_X) \Rightarrow (g \text{ сюръективно}) \land (f \text{ инъективно}).
\]

$\blacktriangleleft$ Действительно, если $f \colon X \to Y$, $g \colon Y \to X$ и $g \circ f = e_X \colon X \to X$, то
\[
  X = e_X(X) = (g \circ f)(X) = g(f(X)) \subset g(Y)
\]
и, значит, $g$ сюръективно.

Далее, если $x_1 \in X$ и $x_2 \in X$, то
\begin{align*}
  (x_1 \neq x_2) &\Rightarrow (e_X(x_1) \neq e_X(x_2)) \Rightarrow ((g \circ f)(x_1) \neq (g \circ f)(x_2)) \Rightarrow \\
  &\Rightarrow (g(f(x_1)) \neq g(f(x_2))) \Rightarrow (f(x_1) \neq f(x_2)),
\end{align*}
следовательно, $f$ инъективно. $\blacktriangleright$

\medskip

Через операцию композиции отображений можно описать взаимно обратные отображения.

\medskip

\textbf{Утверждение.} \textit{Отображения $f \colon X \to Y$, $g \colon Y \to X$ являются биективными и взаимно обратными в том и только в том случае, когда $g \circ f = e_X$ и
$f \circ g = e_Y$.}

$\blacktriangleleft$ В силу леммы одновременное выполнение условий $g \circ f = e_X$ и $f \circ g = e_Y$
гарантирует сюръективность и инъективность, т.\,е. биективность каждого
из отображений $f$, $g$.

Эти же условия показывают, что $y = f(x)$ в том и только в том случае,
когда $x = g(y)$. $\blacktriangleright$

Выше мы исходили из явного построения обратного отображения. Из доказанного утверждения следует, что мы могли бы дать менее наглядное, но
зато более симметричное определение взаимно обратных отображений как
таких, которые удовлетворяют двум условиям: $g \circ f = e_X$ и $f \circ g = e_Y$ (см. в
этой связи задачу 6 в конце параграфа).

\subsection*{4. Функция как отношение. График функции}

В заключение вернемся
вновь к самому понятию функции. Отметим, что оно претерпело длительную и довольно сложную эволюцию.

Термин <<функция>> впервые появился в период 1673--1692 г. у Г.\,Лейбница (правда, в некотором более узком смысле). В смысле, близком к современному, этот термин установился к 1698 г. в переписке Иоганна Бернулли\footnote{И.\,Бернулли (1667--1748) --- один из ранних представителей знаменитого семейства
швейцарских ученых Бернулли; аналитик, геометр, механик. Стоял у истоков вариационного исчисления. Дал первое систематическое изложение дифференциального и интегрального исчисления.} с
Лейбницем.

Описание функции, почти совпадающее с приведенным в начале параграфа, встречается уже у Эйлера (середина XVIII столетия). К началу XIX
века оно появляется уже в учебниках математики С.\,Лакруа,\footnote{С.\,Ф.\,Лакруа (1765--1843) --- французский математик и педагог (профессор Нормальной и Политехнической школ, член Парижской академии наук).} переведенных
на русский язык. Активным сторонником такого понимания функции был
Н.\,И.\,Лобачевский\footnote{Н.\,И.\,Лобачевский (1792--1856) --- великий русский ученый, которому, наряду с великим немецким естествоиспытателем К.\,Ф.\,Гауссом (1777--1855) и выдающимся венгерским математиком Я.\,Бойяи (1802--1860), принадлежит честь открытия неевклидовой
геометрии, носящей его имя.}. Более того, Н.\,И.\,Лобачевский указал, что <<обширный
взгляд теории допускает существование зависимости только в том смысле,
чтобы числа одни с другими в связи понимать как бы \textit{данными вместе}>>\footnote{\textit{Лобачевский Н.\,И.} Полное собр. соч. Т.\,5. М.--Л.: Гостехиздат, 1951. С.\,44.}.
Это и есть идея точного определения понятия функции, которое мы теперь
собираемся изложить.

Приведенное в начале параграфа описание понятия функции представляется весьма динамичным и отражающим суть дела. Однако с точки зрения
современных канонов оно не может быть названо определением, ибо использует эквивалентное функции понятие соответствия. Для сведения читателя мы укажем здесь, каким образом дается определение функции на языке
теории множеств. (Интересно, что понятие отношения, к которому мы сейчас обратимся, и у Лейбница предшествовало понятию функции.)

\medskip

\textbf{a. Отношение.} \textit{Отношением} $\mathscr{R}$ называют любое множество упорядоченных пар $(x, y)$.

Множество $X$ первых элементов упорядоченных пар, составляющих $\mathscr{R}$,
называют \textit{областью определения отношения} $\mathscr{R}$, а множество $Y$ вторых элементов этих пар --- \textit{областью значений} отношения $\mathscr{R}$.

Вместо того чтобы писать $(x, y) \in \mathscr{R}$, часто пишут $x \mathrel{\mathscr{R}} y$ и говорят, что $x$
\textit{связано с $y$ отношением} $\mathscr{R}$.

Если $\mathscr{R} \subset X^2$, то говорят, что \textit{отношение $\mathscr{R}$ задано на} $X$.

Рассмотрим несколько примеров.

\medskip

\textbf{Пример~13.} Диагональ
\[
  \Delta = \{(a, b) \in X^2 \mid a = b\}
\]
есть подмножество $X^2$, задающее отношение равенства между элементами
множества $X$. Действительно, $a \mathrel{\Delta} b$ означает $(a, b) \in \Delta$, т.\,е. $a = b$.

\textbf{Пример~14.} Пусть $X$ --- множество прямых в плоскости.

Две прямые $a \in X$ и $b \in X$ будем считать находящимися в отношении $\mathscr{R}$
и будем писать $a \mathrel{\mathscr{R}} b$, если прямая $b$ параллельна прямой $a$. Ясно, что тем
самым в $X^2$ выделяется множество $\mathscr{R}$ пар $(a, b)$ таких, что $a \mathrel{\mathscr{R}} b$. Из курса
геометрии известно, что отношение параллельности между прямыми обладает следующими свойствами:

$a \mathrel{\mathscr{R}} a$ (рефлексивность);

$a \mathrel{\mathscr{R}} b \Rightarrow b \mathrel{\mathscr{R}} a$ (симметричность);

$(a \mathrel{\mathscr{R}} b) \land (b \mathrel{\mathscr{R}} c) \Rightarrow a \mathrel{\mathscr{R}} c$ (транзитивность).

\medskip

Любое отношение $\mathscr{R}$, обладающее перечисленными тремя свойствами,
т.\,е. рефлексивное\footnote{Полезно для полноты отметить, что отношение $\mathscr{R}$ называется \textit{рефлексивным}, если его
область определения и область значений совпадают и для любого элемента $a$ из области
определения отношения $\mathscr{R}$ выполнено $a \mathrel{\mathscr{R}} a$.}, симметричное и транзитивное, принято называть \textit{отношением эквивалентности}. Отношение эквивалентности обозначается специальным символом $\sim$, который в этом случае ставится вместо буквы $\mathscr{R}$,
обозначающей отношение. Итак, в случае отношения эквивалентности будем писать $a \sim b$ вместо $a \mathrel{\mathscr{R}} b$ и говорить, что $a$ \textit{эквивалентно} $b$.

\textbf{Пример~15.} Пусть $M$ --- некоторое множество, а $X = \mathscr{P}(M)$ --- совокупность всех его подмножеств. Для двух произвольных элементов $a$ и $b$
множества $X = \mathscr{P}(M)$, т.\,е. для двух подмножеств $a$ и $b$ множества $M$, всегда
выполнена одна из следующих трех возможностей: $a$ содержится в $b$; $b$ содержится в $a$; $a$ не является подмножеством $b$ и $b$ не является подмножеством $a$.

Рассмотрим в качестве отношения $\mathscr{R}$ в $X^2$ отношение включения для подмножеств $X$, т.\,е. положим по определению
\[
  a \mathrel{\mathscr{R}} b := (a \subset b).
\]

Это отношение, очевидно, обладает следующими свойствами:

$a \mathrel{\mathscr{R}} a$ (рефлексивность);

$(a \mathrel{\mathscr{R}} b) \land (b \mathrel{\mathscr{R}} c) \Rightarrow a \mathrel{\mathscr{R}} c$ (транзитивность);

$(a \mathrel{\mathscr{R}} b) \land (b \mathrel{\mathscr{R}} a) \Rightarrow a \mathrel{\Delta} b$, т.\,е. $a = b$ (антисимметричность).

\medskip

Отношение между парами элементов некоторого множества $X$, обладающее указанными тремя свойствами, принято называть отношением \textit{частичного порядка} на множестве $X$. Для отношения частичного порядка вместо
$a \mathrel{\mathscr{R}} b$ часто пишут $a \preccurlyeq b$ и говорят, что $b$ \textit{следует за}~$a$.

Если кроме отмеченных двух свойств, определяющих отношение частичного порядка, выполнено условие, что
\[
  \forall a\; \forall b\; ((a \mathrel{\mathscr{R}} b) \lor (b \mathrel{\mathscr{R}} a)),
\]
т.\,е. любые два элемента множества $X$ сравнимы, то отношение $\mathscr{R}$ называется \textit{отношением порядка}, а множество $X$ с определенным на нем отношением
порядка называется \textit{линейно упорядоченным}.

Происхождение этого термина связано с наглядным образом числовой
прямой $\mathbb{R}$, на которой действует отношение $a \leqslant b$ между любой парой вещественных чисел.

\medskip

\textbf{b. Функция и график функции.} Отношение $\mathscr{R}$ называется \textit{функциональным}, если
\[
  (x \mathrel{\mathscr{R}} y_1) \land (x \mathrel{\mathscr{R}} y_2) \Rightarrow (y_1 = y_2).
\]

Функциональное отношение называют \textit{функцией}.

В частности, если $X$ и $Y$ --- два не обязательно различных множества, то
определенное на $X$ отношение $\mathscr{R} \subset X \times Y$ между элементами $x$ из $X$ и $y$ из $Y$
\textit{функционально}, если для любого $x \in X$ существует и притом единственный
элемент $y \in Y$, находящийся с $x$ в рассматриваемом отношении, т.\,е. такой,
для которого $x \mathrel{\mathscr{R}} y$.

Такое функциональное отношение $\mathscr{R} \subset X \times Y$ и есть \textit{отображение из $X$
в $Y$}, или \textit{функция из $X$ в~$Y$}.

Функции мы чаще всего будем обозначать символом $f$. Если $f$ --- функция,
то вместо $x \mathrel{f} y$ мы по-прежнему будем писать $y = f(x)$ или $x \xmapsto{f} y$, называя $y = f(x)$ \textit{значением} функции $f$ на элементе $x$ или \textit{образом} элемента $x$ при
отображении $f$.

Сопоставление по <<закону>> $f$ элементу $x \in X$ <<соответствующего>> элемента $y \in Y$, о чем говорилось в исходном описании понятия функции, как видим, состоит в том, что для каждого $x \in X$ указывается тот единственный
элемент $y \in Y$, что $x \mathrel{f} y$, т.\,е. $(x, y) \in f \subset X \times Y$.

Графиком функции $f \colon X \to Y$, понимаемой в смысле исходного описания,
называют подмножество $\Gamma$ прямого произведения $X \times Y$, элементы которого
имеют вид $(x, f(x))$. Итак,
\[
  \Gamma := \{(x, y) \in X \times Y \mid y = f(x)\}.
\]

В новом описании понятия функции, когда мы ее задаем как подмножество $f \subset X \times Y$, конечно, уже нет разницы между функцией и ее графиком.

Мы указали на принципиальную возможность формального теоретикомножественного определения функции, сводящуюся по существу к отождествлению функции и ее графика. Однако мы не собираемся в дальнейшем
ограничиваться только такой формой задания функции. Функциональное
отношение иногда удобно задать в аналитической форме, иногда таблицей
значений, иногда словесным описанием процесса (алгоритма), позволяющего по данному $x \in X$ находить соответствующий элемент $y \in Y$. При
каждом таком способе задания функции имеет смысл вопрос о ее задании
с помощью графика, что формулируют так: построить график функции.
Задание числовых функций хорошим графическим изображением часто
бывает полезно тем, что делает наглядным основные качественные особенности функциональной зависимости. Для расчетов графики тоже можно
использовать (номограммы), но, как правило, в тех случаях, когда расчет не
требует высокой точности. Для точных расчетов используют табличное задание функции, а чаще --- алгоритмическое, реализуемое в вычислительных
машинах.

\subsection*{Упражнения}

\textbf{1.} Композиция $\mathscr{R}_2 \circ \mathscr{R}_1$ отношений $\mathscr{R}_1$, $\mathscr{R}_2$ определяется следующим образом:
\[
  \mathscr{R}_2 \circ \mathscr{R}_1 := \{(x, z) \mid \exists y\; (x \mathrel{\mathscr{R}_1} y) \land (y \mathrel{\mathscr{R}_2} z)\}.
\]

В частности, если $\mathscr{R}_1 \subset X \times Y$ и $\mathscr{R}_2 \subset Y \times Z$, то $\mathscr{R} = \mathscr{R}_2 \circ \mathscr{R}_1 \subset X \times Z$, причем
\[
  x \mathrel{\mathscr{R}} z := \exists y\; ((y \in Y) \land (x \mathrel{\mathscr{R}_1} y) \land (y \mathrel{\mathscr{R}_2} z)).
\]

a) Пусть $\Delta_X$ --- диагональ множества $X^2$, а $\Delta_Y$ --- диагональ множества $Y^2$.
Покажите, что если отношения $\mathscr{R}_1 \subset X \times Y$ и $\mathscr{R}_2 \subset Y \times X$ таковы, что $(\mathscr{R}_2 \circ \mathscr{R}_1 = \Delta_X) \land
{} \land (\mathscr{R}_1 \circ \mathscr{R}_2 = \Delta_Y)$, то оба они функциональны и задают взаимно обратные отображения множеств $X$,~$Y$.

b) Пусть $\mathscr{R} \subset X^2$. Покажите, что условие транзитивности отношения $\mathscr{R}$ равносильно тому, что $\mathscr{R} \circ \mathscr{R} \subset \mathscr{R}$.

c) Отношение $\mathscr{R}' \subset Y \times X$ называется \textit{транспонированным} отношением $\mathscr{R} \subset X \times Y$,
если $(y \mathrel{\mathscr{R}'} x) \Leftrightarrow (x \mathrel{\mathscr{R}} y)$.

Покажите, что антисимметричность отношения $\mathscr{R} \subset X^2$ равносильна условию
$\mathscr{R} \cap \mathscr{R}' \subset \Delta_X$.

d) Проверьте, что любые два элемента множества $X$ связаны (в том или ином
порядке) отношением $\mathscr{R} \subset X^2$, если и только если $\mathscr{R} \cup \mathscr{R}' = X^2$.

\medskip

\textbf{2.} Пусть $f \colon X \to Y$ --- отображение. Прообраз $f^{-1}(y) \subset X$ элемента $y \in Y$ называется
\textit{слоем} над~$y$.

a) Укажите слои для отображений
\[
  \operatorname{pr}_1 \colon X_1 \times X_2 \to X_1, \quad \operatorname{pr}_2 \colon X_1 \times X_2 \to X_2.
\]

b) Элемент $x_1 \in X$ будем считать связанным с элементом $x_2 \in X$ отношением $\mathscr{R} \subset
{} \subset X^2$ и писать $x_1 \mathrel{\mathscr{R}} x_2$, если $f(x_1) = f(x_2)$, т.\,е. если $x_1$ и $x_2$ лежат в одном слое.

Проверьте, что $\mathscr{R}$ есть отношение эквивалентности.

c) Покажите, что слои отображения $f \colon X \to Y$ не пересекаются, а объединением
слоев является все множество $X$.

d) Проверьте, что любое отношение эквивалентности между элементами множества позволяет представить это множество в виде объединения непересекающихся
классов эквивалентных элементов.

\medskip

\textbf{3.} Пусть $f \colon X \to Y$ --- отображение из $X$ в $Y$. Покажите, что если $A$ и $B$ --- подмножества $X$, то

a) $(A \subset B) \Rightarrow (f(A) \subset f(B)) \not\Rightarrow (A \subset B)$, \quad b) $(A \neq \varnothing) \Rightarrow (f(A) \neq \varnothing)$,

c) $f(A \cap B) \subset f(A) \cap f(B)$, \qquad d) $f(A \cup B) = f(A) \cup f(B)$;

если $A'$ и $B'$ --- подмножества $Y$, то

e) $(A' \subset B') \Rightarrow (f^{-1}(A') \subset f^{-1}(B'))$, \qquad f) $f^{-1}(A' \cap B') = f^{-1}(A') \cap f^{-1}(B')$,

g) $f^{-1}(A' \cup B') = f^{-1}(A') \cup f^{-1}(B')$;

если $Y \supset A' \supset B'$, то

h) $f^{-1}(A' \setminus B') = f^{-1}(A') \setminus f^{-1}(B')$, \qquad i) $f^{-1}(C_Y A') = C_X f^{-1}(A')$;

для любого множества $A \subset X$ и любого множества $B' \subset Y$

j) $f^{-1}(f(A)) \supset A$, \qquad k) $f(f^{-1}(B')) \subset B'$.

\medskip

\textbf{4.} Покажите, что отображение $f \colon X \to Y$

a) сюръективно, если и только если для любого множества $B' \subset Y$ справедливо
$f(f^{-1}(B')) = B'$;

b) биективно, если и только если для любого множества $A \subset X$ и любого множества $B' \subset Y$ справедливо
\[
  (f^{-1}(f(A)) = A) \land (f(f^{-1}(B')) = B').
\]

\medskip

\textbf{5.} Проверьте эквивалентность следующих утверждений относительно отображения $f \colon X \to Y$:

a) $f$ инъективно;

b) $f^{-1}(f(A)) = A$ для любого множества $A \subset X$;

c) $f(A \cap B) = f(A) \cap f(B)$ для любой пары $A, B$ подмножеств $X$;

d) $f(A) \cap f(B) = \varnothing \Leftrightarrow A \cap B = \varnothing$;

e) $f(A \setminus B) = f(A) \setminus f(B)$, если $X \supset A \supset B$.

\medskip

\textbf{6.} a) Если отображения $f \colon X \to Y$ и $g \colon Y \to X$ таковы, что $g \circ f = e_X$, где $e_X$ --- тождественное отображение множества $X$, то $g$ называется \textit{левым обратным отображением} для $f$, а $f$ --- \textit{правым обратным отображением} для $g$. Покажите, что, в отличие от единственного
обратного отображения, может существовать много односторонних обратных отображений.

Рассмотрите, например, отображения $f \colon X \to Y$ и $g \colon Y \to X$, где $X$ --- одноэлементное, а $Y$ --- двухэлементное множества, или отображения последовательностей
\begin{align*}
  (x_1, \ldots, x_n, \ldots\,) &\xrightarrow{f_n} (a, x_1, \ldots, x_n, \ldots\,), \\
  (y_2, \ldots, y_n, \ldots\,) &\xleftarrow{g} (y_1, y_2, \ldots, y_n, \ldots\,).
\end{align*}

b) Пусть $f \colon X \to Y$ и $g \colon Y \to Z$ --- биективные отображения. Покажите, что отображение $g \circ f \colon X \to Z$ биективно и что $(g \circ f)^{-1} = f^{-1} \circ g^{-1}$.

c) Покажите, что для любых отображений $f \colon X \to Y$, $g \colon Y \to Z$ и любого множества $C \subset Z$ справедливо равенство
\[
  (g \circ f)^{-1}(C) = f^{-1}(g^{-1}(C)).
\]

d) Проверьте, что отображение $F \colon X \times Y \to Y \times X$, задаваемое соответствием
$(x, y) \mapsto (y, x)$, биективно. Опишите взаимосвязь графиков взаимно обратных отображений $f \colon X \to Y$ и $f^{-1} \colon Y \to X$.

\medskip

\textbf{7.} a) Покажите, что при любом отображении $f \colon X \to Y$ определяемое соответствием $x \xmapsto{f} (x, f(x))$, является инъективным.

b) Пусть частица движется равномерно по окружности $Y$; пусть $X$ --- ось времени и $x \xmapsto{f} y$ --- соответствие между моментом времени $x \in X$ и положением $y = f(x) \in Y$
частицы. Изобразите график функции $f \colon X \to Y$ в $X \times Y$.

\medskip

\textbf{8.} a) Для каждого из разобранных в \S\,3 примеров 1--12 выясните, является ли
указанное в нем отображение сюръективным, инъективным, биективным или оно
не принадлежит ни одному из указанных классов.

b) Закон Ома $I = V / R$ связывает силу тока $I$ в проводнике с напряжением $V$ на
концах проводника и сопротивлением $R$ проводника. Укажите, отображение $O \colon X \to
{} \to Y$ каких множеств соответствует закону Ома. Подмножеством какого множества
является отношение, отвечающее закону Ома?

c) Найдите преобразования $G^{-1}$, $L^{-1}$, обратные к преобразованиям Галилея и Лоренца.

\medskip

\textbf{9.} a) Множество $S \subset X$ называется \textit{устойчивым} относительно отображения
$f \colon X \to X$, если $f(S) \subset S$. Опишите множества, устойчивые относительно сдвига
плоскости на данный лежащий в ней вектор.

b) Множество $I \subset X$ называется \textit{инвариантным} относительно отображения
$f \colon X \to X$, если $f(I) = I$. Опишите множества, инвариантные относительно поворота
плоскости вокруг фиксированной точки.

c) Точка $p \in X$ называется \textit{неподвижной точкой} отображения $f \colon X \to X$, если
$f(p) = p$. Проверьте, что любая композиция сдвига, вращения и гомотетии плоскости
имеет неподвижную точку, если коэффициент гомотетии меньше единицы.

d) Считая преобразования Галилея и преобразования Лоренца отображениями
плоскости на себя, при которых точка с координатами $(x, t)$ переходит в точку с координатами $(x', t')$, найдите инвариантные множества этих преобразований.

\medskip

\textbf{10.} Рассмотрим установившийся поток жидкости (т.\,е. скорость в каждой точке
потока не меняется со временем). За время $t$ частица, находящаяся в точке $x$ потока,
переместится в некоторую новую точку $f_t(x)$ пространства. Возникающее отображение $x \mapsto f_t(x)$ точек пространства, занимаемого потоком, зависит от времени $t$
и называется \textit{преобразованием за время}~$t$. Покажите, что $f_{t_2} \circ f_{t_1} = f_{t_1} \circ f_{t_2} = f_{t_1 + t_2}$ и
$f_t \circ f_{-t} = f_0 = e_X$.
